% =============================================================================
%  04_category.tex —— §4 范畴论
% =============================================================================
\section{范畴论}\label{sec:cat}

一般地，我们希望把\emph{对象}（层、向量丛、阿贝尔群、模等等）换成由对象组成的
复形，而这些复形对某个特定的\emph{运算}来说\emph{“性质更好”}，例如对
$\otimes A$、$\Hom(\,\cdot\,,A)$、$\Hom(A,\,\cdot\,)$、$\Gamma(\,\cdot\,)$
等等（“性质更好”的意思是：把运算作用到这类对象上时\emph{不丢信息}）。
有时这样的复形是典范的；一般情形下，我们希望把它的选择做得\emph{自然}。

上面这些斜体的词是在暗示范畴论：阿贝尔范畴是这套一般理论的天然舞台，前面的词
分别换成范畴的对象、无环对象、函子、正合、函子化。下面就来定义这些东西。见了
“范畴”二字就发怵的读者，不妨把它想成复流形上的（层构成的）范畴。

\begin{definition}\label{def:additive}
\term{加性范畴}{additive category}是一个范畴 $\A$，满足
\begin{itemize}
\item 每个态射集 $\Hom(A,B)$ 构成阿贝尔群；
\item 态射的复合对由该阿贝尔群结构给出的加法分配，即
  $f\circ(g+h)=f\circ g+f\circ h$ 与 $(f+g)\circ h=f\circ h+g\circ h$；
\item 任意两个对象 $A$、$B$ 都有乘积（直和）$A\times B$，满足通常的泛性质
  （例如见 \cite{Ma}）；
\item 存在\term{零对象}{zero object} $0$，使 $\Hom(0,0)$ 是零群
  （即只含恒等态射）。于是对一切 $A$ 有 $\Hom(0,A)=0=\Hom(A,0)$，
  而任意两个对象之间唯一的零态射就是经过零对象分解的那个。
\end{itemize}
\term{阿贝尔范畴}{abelian category}是还满足下述条件的加性范畴：
\begin{itemize}
\item 一切态射都有核与余核，并且单态射（核为零的态射）是其余核映射的核，
  满态射（余核为零的态射）是其核的余核。
\end{itemize}
\end{definition}

上面提到的核与余核，在一般范畴中由显然的泛性质定义（例如 $f\colon A\to B$ 的核
是一个 $K\to A$，使得 $g\colon C\to A$ 当且仅当 $f\circ g=0$ 时唯一地经过它
分解）。这类事情请参看 \cite{Ma}；我们只关心具体范畴（对象是集合的范畴），
此时核与余核就是通常的核与余核，最后一条公理自动成立。

于是在阿贝尔范畴中可以谈论正合列与\emph{链复形}，以及复形的上同调。阿贝尔范畴
之间的加性函子是\emph{正合}的（相应地，左正合或右正合），如果它保持正合列
（相应地，保持短正合列 $0\to A\to B\to C$ 或 $A\to B\to C\to 0$）。

\begin{definition}\label{def:db}
阿贝尔范畴 $\A$ 的\term{有界导出范畴}{bounded derived category} $D^b(\A)$
以有界（即长度为有限）的 $\A$-链复形为对象，态射则由链映射给出，并把拟同构
如下形式地可逆化（\cite{GM} III 2.2）：对每个链映射 $f\colon X_f\to Y_f$ 引入
态射 $f$，对每个拟同构 $g\colon X_g\xrightarrow{\sim}Y_g$ 引入态射
$g^{-1}\colon Y_g\to X_g$；然后取这些态射的一切乘积，要求前一个的余定义域是
后一个的定义域；最后把任何组合 $f_1f_2$ 与复合 $f_1\circ f_2$ 等同起来，
把 $gg^{-1}$、$g^{-1}g$ 与相应的恒等态射 $\id_{Y_g}$、$\id_{X_g}$ 等同起来。
\end{definition}

\begin{remark}\label{rem:four}
同样可以定义无界导出范畴，以及有界下方、有界上方复形的范畴 $D^+(\A)$、
$D^-(\A)$。下文用 $D(\A)$ 泛指这四种导出范畴中的一种。
\end{remark}

$D(\A)$ 中的态射由\term{屋顶}{roof}表示：
$$
\begin{array}{ccc}
& Z\udot & \\
{\scriptstyle s}\swarrow & & \searrow{\scriptstyle g} \\
X\udot & \dashrightarrow{\scriptstyle f} & Y\udot,
\end{array}
$$
其中 $s$ 是拟同构，并令 $f=gs^{-1}$。（任何态射显然都是这种屋顶的复合；一个
屋顶就够用，见 \cite{GM} III 2.8。）把拟同构可逆化的这套“\term{局部化}{localisation}”
手续有一些值得注意的性质（例如我们会看到，同伦的映射在导出范畴中互相等同，
给出同一个态射）。这些性质使 $D(\A)$ 成为加性范畴，但我们也将看到它没有核与
余核。还是老办法：回到拓扑，看为什么没有、它们被什么取代，以及 $D(\A)$ 的结构
究竟是什么（既然它不是阿贝尔范畴的结构）。

\begin{conceptbox}{导出范畴的“三步走”}
从 $\A$ 走到 $D^b(\A)$，可以记成三步：
\begin{enumerate}[leftmargin=1.6em,itemsep=2pt,topsep=2pt]
\item \textbf{只看复形}：对象换成有界复形，忘记“单个对象”的执念；
\item \textbf{把拟同构当同构}：形式地添入拟同构的逆，得到局部化范畴；
\item \textbf{接受非阿贝尔}：核与余核消失，取而代之的是锥与特异三角（§\ref{sec:cone}）。
\end{enumerate}
第 2 步的严格实现有两套等价说法：一是如定义 \ref{def:db} 那样自由添入逆元再作
商；二是用“屋顶”$X\xleftarrow{s}Z\xrightarrow{g}Y$ 作态射，并规定两个屋顶等价
当且仅当存在第三个屋顶同时连通它们。后者正是分数演算（calculus of fractions），
它成立的关键是拟同构对“拉回”封闭。
[来源：Gelfand--Manin 第 III 章 §2；Weibel 第 10 章；A 级]
\end{conceptbox}

\begin{yizhu}[为什么“单态射=其余核的核”是阿贝尔范畴的关键]
这条公理把“单/满”与“核/余核”这两套语言接了起来，从而保证
$\operatorname{im}f\cong\operatorname{coim}f$，正合列才可以用核与余核刻画。
在一般加性范畴中，单态射不一定是某个映射的核；阿贝尔范畴要求这点，正是为了
让同调代数（蛇引理、五引理、长正合列）能照常运转。
[来源：MacLane《Categories for the Working Mathematician》第 VIII 章；A 级]
\end{yizhu}

\begin{yizhu}[定义 \ref{def:db} 为何只添“一步”就够]
定义里只形式地添入拟同构的逆，再取一切乘积并作商。这样得到的态射集合，与
“把拟同构全体可逆化”的普遍性质所要求的局部化，恰好一致——因为拟同构构成一个
\emph{饱和系}，满足分数演算的两条条件（右 Ore 条件与消去条件）。因此每个态射都
可以写成 $g s^{-1}$ 的形状，不必添入无穷长的词。
[来源：Gelfand--Manin III 2.8；Kashiwara--Schapira《Sheaves on Manifolds》第 10 章；A 级]
\end{yizhu}

\begin{keytakeaways}
  \item 加性范畴给 $\Hom$ 群结构；阿贝尔范畴再加“单态射 = 其余核的核”，
        于是有核、余核与正合列。
  \item $D^b(\A)$ 的对象是有界复形，态射是“链映射 + 拟同构的逆”的形式复合；
        屋顶 $X\xleftarrow{s}Z\xrightarrow{g}Y$ 是它的标准代表。
  \item 只添一步逆就够，因为拟同构类满足 Ore 条件（译注 14）。
  \item $D^b(\A)$ 不是阿贝尔范畴：核与余核不存在，取而代之的是锥与特异三角（§5）。
\end{keytakeaways}
